\section{Inbetriebnahme}
\label{sec:start}

\subsection{Einschalten}
\begin{steps}
    \item Alle Prüfleitungen von den Buchsen entfernen.
    \item \tbc{Ein-Taste} drücken. Der Startbildschirm mit dem \ST{}-Logo erscheint.
    \item Nach dem Start wird das Hauptmenü angezeigt (\autoref{fig:idle}).
\end{steps}
\todo[inline]{Genauer Einschaltvorgang (Taste, Haltezeit) für die finale Hardware}

\begin{note}
    Liegt an den Buchsen eine Versorgung mit \qtyrange{15}{24}{\volt} an, beginnt das Gerät
    automatisch zu laden (\autoref{sec:charging}).
\end{note}

\subsection{Betriebsart wählen}
\begin{steps}
    \item \enc{} drehen, um eine Betriebsart zu wählen. Ihr Name wird unter dem Symbol angezeigt.
    \item \btn{OK} drücken, um die Betriebsart zu öffnen.
    \item \btn{ESC} drücken, um ins Hauptmenü zurückzukehren. Der Ausgang wird ausgeschaltet.
\end{steps}

\subsection{Ausschalten}
\begin{itemize}
    \item Im Hauptmenü \btn{ESC} drücken (\disp{ESC: Shutdown}).
    \item Das Gerät schaltet nach \qty{10}{\minute} ohne Bedienung (\enc{} drehen oder Taste drücken)
          automatisch aus. Gehaltenes \btn{OUT} zählt als Bedienung. Beim Laden gilt dies nicht.
\end{itemize}

\subsection{Batterie}
Das Gerät wird von einer internen \ac{LiFePO4}-Batterie versorgt
(4 Zellen, \qty{12.8}{\volt}, \tbc{\qty{4.5}{\ampere\hour}}).
\todo[inline]{Wo wird der Ladezustand im normalen Betrieb angezeigt? (Anforderung R:Bat\_Status)}
\todo[inline]{Erwartete Laufzeit je Betriebsart}
